\documentclass[10pt,a4paper]{amsart}
\usepackage[margin=19mm]{geometry}
\usepackage{amsmath,amssymb,mathtools}
\usepackage[scr=rsfs,cal=euler]{mathalfa}
\usepackage{xfrac}
\usepackage[unicode,hidelinks,bookmarks=false]{hyperref}
\newcommand{\ie}{i.e.\ }
\newcommand{\cf}{cf.\ }
\newcommand{\resp}{respectively\ }
\newcommand{\Subsection}{\subsection}
\providecommand{\nobreakdash}{\nobreak\hbox{-}}
\setlength{\emergencystretch}{2em}
\multlinegap0pt
\usepackage{amssymb}
\usepackage{mathtools}
\usepackage{tikz}
\usepackage{float}
\usepackage[shortlabels]{enumitem}
\newtheorem{theorem}{Theorem}
\title{Metric Bases of Barycentric and Matching Subdivisions of Zero-Divisor Graphs}
\author{Vidya S, Prasanna Poojary, Bilal Ahmad Rather, Vadiraja Bhatta G R}
\date{Source: arXiv:2609.09896v1 (CC BY 4.0)}
\begin{document}
\maketitle

\noindent\emph{Source and licence.} Metric Bases of Barycentric and Matching Subdivisions of Zero-Divisor Graphs. Version arXiv:2609.09896v1 (CC BY 4.0), distributed by arXiv under the Creative Commons Attribution 4.0 International licence, http://creativecommons.org/licenses/by/4.0/, as recorded in the arXiv licence metadata for this exact version and shown on the abstract page https://arxiv.org/abs/2609.09896v1. Full licence notice: \texttt{licenses/arxiv-2609.09896-CC-BY-4.0.txt}.\par\smallskip

\noindent\textit{Editorial scope.} Counted unit: the article's theorem thm-general-matching (source0.tex lines 547--552). The article's complete original proof of it is delivered with it (source0.tex lines 554--602, one \texttt{proof} environment of 49 lines). No mathematics is written, repaired or re-derived: the statement and the proof are the article's own lines, in the article's own notation, and no retained line is substituted or re-flowed.  No further source-local context is needed: every cross-reference of every retained line resolves inside the two delivered spans.  Nothing was omitted from any retained span, so the package carries no omission record.  No retained line contains a citation, so the wrapper carries no bibliography and no bibliography entry.  No retained line contains a figure environment, an image-inclusion command or drawing code, so no image is delivered and the package declares no asset dependency.  Editorial formatting only, in addition: an AMS \texttt{amsart} preamble that restates the article's own declarations and macros used by the retained lines, namely the environments \texttt{theorem} on separate counters, together with the standard packages those lines need.  The retained environments print in this document as Theorem 1 in order of appearance, whereas the article prints the same items with the numbering of its own full document.  The complete native source of the article as distributed in the arXiv e-print for this exact version is bundled unchanged as \texttt{sources/209859/source0.tex}.
\begin{theorem}\label{thm-general-matching}
Let $G_r$ be obtained from $\Gamma(\mathbb Z_{pq})$ by subdividing the edges of a matching of size $r$, where $0\leq r\leq p-2$. Then
$$
\dim(G_r)=p+q-r-4.
$$
\end{theorem}

\begin{proof}
Write $a=p-1$ and $b=q-1$, so $\Gamma(\mathbb Z_{pq})\cong K_{a,b}$. Without loss of generality, suppose that the subdivided matching edges are $q_ip_i$ for $1\leq i\leq r$, and let $m_i$ be the subdivision vertex on $q_ip_i$.

First we construct a resolving set. Define
$$
W=\{q_1,q_2,\ldots,q_{a-1}\}\cup\{p_{r+1},p_{r+2},\ldots,p_{b-1}\}.
$$
The first set has $a-1=p-2$ vertices and the second has $b-r-1=q-r-2$ vertices. Therefore
$$
|W|=p+q-r-4.
$$
We verify that $W$ resolves $G_r$. Two vertices of $S_q$ are separated by a selected $q$-vertex unless one of them is the omitted vertex $q_a$; in that remaining case, a selected vertex $p_{r+1}$, when needed, or one of the selected $q$-vertices separates the pair by the usual $0$ and $2$ coordinates. The vertices $p_{r+1},\ldots,p_{b-1}$ similarly separate all vertices of $S_p$ except the single omitted vertex $p_b$, and any matched vertex $p_i$ is separated from $p_b$ by $q_i$, since $d(q_i,p_i)=2$ and $d(q_i,p_b)=1$. A pair consisting of one vertex from $S_q$ and one from $S_p$ is separated either by a selected vertex of $S_q$ or by a selected vertex of $S_p$; the subdivided pair $q_i,p_i$ is separated by $q_i$, because the corresponding distances are $0$ and $2$. Finally, $m_i$ is separated from $q_i$ and $p_i$ by $q_i$, and if $m_i\neq m_j$, then $q_i$ separates them because $d(q_i,m_i)=1$ while $d(q_i,m_j)=2$. Hence $W$ is resolving, and
$$
\dim(G_r)\leq p+q-r-4.
$$

For the lower bound, let $R$ be any resolving set. The unmatched vertices
$$
Q_0=\{q_{r+1},q_{r+2},\ldots,q_a\}
$$
form a twin class, because they are adjacent to exactly the same vertices of $S_p$ and to no subdivision vertex. Therefore
$$
|R\cap Q_0|\geq |Q_0|-1=a-r-1.
$$
Similarly,
$$
P_0=\{p_{r+1},p_{r+2},\ldots,p_b\}
$$
is a twin class, so
$$
|R\cap P_0|\geq |P_0|-1=b-r-1.
$$
For $1\leq i\leq r$, put $C_i=\{q_i,p_i,m_i\}$. If $R\cap C_i=\emptyset$ and some vertex $q_0\in Q_0$ is not in $R$, then $q_i$ and $q_0$ have the same distance to every landmark outside $C_i$: they are both at distance $1$ from all vertices of $P_0$, both at distance $2$ from all vertices of $Q_0\setminus\{q_0\}$, and have equal distances to all other triples $C_j$ with $j\neq i$. This contradicts the resolving property. Hence, whenever a triple $C_i$ is missed, all vertices of $Q_0$ must be contained in $R$. The same argument with the two partite sets interchanged shows that a missed $C_i$ also forces all vertices of $P_0$ to be contained in $R$.

Moreover, two different triples cannot both be missed. Indeed, if $R\cap C_i=R\cap C_j=\emptyset$ with $i\neq j$, then the vertices $q_i$ and $q_j$ have the same distances to every landmark outside $C_i\cup C_j$, while there are no landmarks inside those two triples. Thus $R$ would not resolve $q_i$ and $q_j$.

Therefore at least $r-1$ of the triples $C_i$ meet $R$. If all $r$ triples meet $R$, then together with the two twin-class bounds we get
$$
|R|\geq (a-r-1)+(b-r-1)+r=a+b-r-2.
$$
If exactly one triple is missed, then all vertices of $Q_0$ and all vertices of $P_0$ are contained in $R$, giving two additional vertices beyond the twin-class lower bounds, and hence
$$
|R|\geq (a-r-1)+(b-r-1)+(r-1)+2=a+b-r-1>a+b-r-2.
$$
In all cases $|R|\geq a+b-r-2=p+q-r-4$. Combining this with the upper bound proves
$$
\dim(G_r)=p+q-r-4.
$$
\end{proof}

\end{document}
